\clearpage
\VerdierPhysicalPageAnchor{0108}%
\begingroup
\small
\setlength{\abovedisplayskip}{4pt plus 2pt minus 2pt}%
\setlength{\belowdisplayskip}{4pt plus 2pt minus 2pt}%

\VerdierSemanticAnchor{ch1-app-signs}
\section*{附录：关于符号选择的说明}
\addcontentsline{toc}{section}{附录：关于符号选择的说明}

\VerdierSemanticAnchor{ch1-app-signs-cone}
映射锥函子由其伴随性质 (3.1.3) 在唯一同构的意义下所确定。为了便于把
$\operatorname{tr}_1(C)$ 的一个态射关联到 $\operatorname{comp}(C)$ 的一个三角，
我们选择“态射 $X\longrightarrow 0$ 的映射锥”作为平移函子（(3.1.7) 和
(3.2.2.7)），也就是说，平移次数并改变微分的符号。

\VerdierSemanticAnchor{ch1-app-signs-strict}
设 $F\colon\prod_{i\in[n]}C_i\longrightarrow C$ 是一个多加法函子。范畴
$\prod_{i\in[n]}\operatorname{comp}(C_i)$ 自然带有一个 $\mathbb Z[n]$-范畴结构。
如果要求函子
\[
  F^{[n]}\colon\prod_{i\in[n]}\operatorname{comp}(C_i)
  \longrightarrow\operatorname{comp}^{n}(C)
\]
是一个严格 $\mathbb Z[n]$-函子，就会得到我们所采用的平移约定：沿一个方向
平移次数，并改变相应微分的符号。

\VerdierSemanticAnchor{ch1-app-signs-simple-complex}
关于关联单复形函子的选择，我们采用了 [1] 的约定（但在我们的多重复形中，
各微分彼此交换）。这使得可以定义函子
$K_{\Pi}F\colon\prod_i K(C_i)\longrightarrow K(C)$。此时
$K_{\Pi}F(X_1,\ldots,X_i[1],\ldots,X_n)$ 与
$K_{\Pi}F(X_1,\ldots,X_i,\ldots,X_n)[1]$ 不同。为了仍能断言函子
$K_{\Pi}F$ 把特异三角送到特异三角，我们必须定义同构

\VerdierSemanticAnchor{ch1-app-signs-shift-isomorphism}
\[
  p_i\colon K_{\Pi}F(X_1,\ldots,X_i[1],\ldots,X_n)
  \xrightarrow{\ \sim\ }
  K_{\Pi}F(X_1,\ldots,X_i,\ldots,X_n)[1]
\]
（见 (3.4.3) 的证明）。

\VerdierSemanticAnchor{ch1-app-signs-anticommutativity}
于是，当 $i\ne j$ 时，得到反交换图：

\VerdierSemanticAnchor{ch1-app-signs-diagram}
\begin{center}
  \resizebox{\textwidth}{!}{%
    $\begin{tikzcd}[ampersand replacement=\&,column sep=3.2em,row sep=4.3em]
      K_{\Pi}F(X_1,\ldots,X_i[1],\ldots,X_j[1],\ldots)
        \arrow[r,"p_j"]
        \arrow[d,"p_i"']
        \& K_{\Pi}F(X_1,\ldots,X_i[1],\ldots,X_j,\ldots)[1]
          \arrow[d,"p_i"] \\
      K_{\Pi}F(X_1,\ldots,X_i,\ldots,X_j[1],\ldots)[1]
        \arrow[r,"p_j"]
        \& K_{\Pi}F(X_1,\ldots,X_i,\ldots,X_j,\ldots)[2].
    \end{tikzcd}$%
  }
\end{center}

\VerdierSemanticAnchor{ch1-app-signs-koszul-choice}
随后我们在 (1.5.2) 中说明了如何把这些同构相容地扩张为同构 $p(\sigma)$，
它对每个 $\sigma\in\mathbb Z[n]$ 都有定义。这里包含一个任意选择：需要在
一个本身已确定的上同调类，即 Koszul 类 (1.6.3) 中选择一个余循环。为了符合
通常接受的规则（例如 [1]），我们在所有情形下都选择 Koszul 余循环 (1.6.6)。
\endgroup
